% !TEX program = xelatex
% 普西函数与黎曼泽塔函数
\documentclass[fontset=windows, chinese, 10pt, xcolor=dvipsnames]{beamer}
% Shared preamble for the chapter 11 standalone lectures.
% Each lecture sets \title and \subtitle after this file is read.

\usepackage{amsmath,amssymb,mathtools}
\usepackage{xeCJK}
\usepackage{zhnumber}
\usepackage{tikz}
\usepackage{tikz-3dplot}
\usetikzlibrary{arrows.meta,calc,positioning}

\usetheme{metropolis}
\usepackage{beamercolorthemeHistNanjing}

\newcommand{\red}[1]{\textcolor{red}{#1}}
\newcommand{\blue}[1]{\textcolor{blue}{#1}}
\newcommand{\green}[1]{\textcolor{green}{#1}}
\newcommand{\gray}[1]{\textcolor{gray}{#1}}
\newcommand{\orange}[1]{\textcolor{orange}{#1}}

\newcommand{\sssub}[1]{ {\mathrm{#1} }}  % script size subscript\\\
\newcommand{\funcd}[2]{ \frac{d{#1}}{d{#2}}} %function full derivative
\newcommand{\funcpd}[2]{ \frac{\partial{#1}}{\partial{#2}}} %function partial derivative
\newcommand{\fnald}[2] { \frac{\delta {#1}}{\delta {#2}}} % functional derivative

\renewcommand{\Re}{\operatorname{Re}}
\renewcommand{\Im}{\operatorname{Im}}
\def\ext{\sssub{ext}}
\def\x{\sssub{x}}
\def\xc{\sssub{xc}}
\def\H{\sssub{H}}
\def\Hx{\sssub{Hx}}
\def\Hxc{\sssub{Hxc}}
\def\c{\sssub{c}}
\def\s{\sssub{s}}

% \def\ext{\sssub{ext}}
% \def\ext{\sssub{ext}}
% \def\ext{\sssub{ext}}
% \def\ext{\sssub{ext}}
% \defsssub{ext}


% definitions of bold faced symbols
\def\brpppp{{\mathbf{r}^{\prime\prime\prime\prime}}}
\def\brppp{{\mathbf{r}^{\prime\prime\prime}}}
\def\brpp{{\mathbf{r}^{\prime\prime}}}
\def\brp{{\mathbf{r}^{\prime}}}
\def\bzp{{\mathbf{z}^{\prime}}}
\def\bxp{{\mathbf{x}^{\prime}}}
\def\tp{{{t}^{\prime}}}
\def\tpp{{{t}^{\prime\prime}}}
\def\tppp{{{t}^{\prime\prime\prime}}}

\def\tbr{{\tilde{\mathbf{r}}}}
\def\bk{{\mathbf{k}}}
\def\br{{\mathbf{r}}}
\def\bp{{\mathbf{p}}}
\def\bq{{\mathbf{q}}}
\def\bv{{\mathbf{v}}}
\def\bz{{\mathbf{z}}}
\def\mbf{{\mathbf{f}}}

\def\bx{{\mathbf{x}}}
\def\bR{{\mathbf{R}}}
\def\bM{{\mathbf{M}}}
\def\bP{{\mathbf{P}}}
\def\bT{{\mathbf{T}}}
\def\bK{{\mathbf{K}}}
\def\bA{{\mathbf{A}}}
\def\bB{{\mathbf{B}}}
\def\bD{{\mathbf{D}}}
\def\bE{{\mathbf{E}}}
\def\bF{{\mathbf{F}}}
\def\bG{{\mathbf{G}}}
\def\bH{{\mathbf{H}}}
\def\bI{{\mathbf{I}}}
\def\bJ{{\mathbf{J}}}
\def\bP{{\mathbf{P}}}
\def\bS{{\mathbf{S}}}
\def\bV{{\mathbf{V}}}
\def\bX{{\mathbf{X}}}
\def\bY{{\mathbf{Y}}}
\def\dd{{\mathop{}\!\mathrm{d}}}
\def\dr{{\!\!\dd\br\,}}
\def\dx{{\!\!\dd\bx\,}}
\def\dy{{\!\!\dd\by\,}}
\def\dk{{\!\!\dd\bk\,}}
\def\dG{{\!\!\dd\bG\,}}


\def\half{{\frac{1}{2}}}
\def\tr{{\mathrm{tr}}}


\def\Arg{{\mathrm{Arg}\;}}
\def\arg{{\mathrm{arg}\;}}
\def\Res{{\mathrm{Res}\;}}

% unit-vector basis (bold under XeLaTeX; \mathbf can look non-bold)
\newcommand{\ue}{\boldsymbol{e}}
\newcommand{\un}{\hat{\boldsymbol{n}}}


\author{罗凯}
\institute{南京理工大学物理学院}
\date{\zhdigits{\the\year},\zhganzhinian{\the\year}年秋}

\usebackgroundtemplate{%
\begin{tikzpicture}[remember picture,overlay]
  \node[opacity=0.05, at=(current page.center)] {
    \includegraphics[width=0.6\paperwidth]{njust-logo.jpeg}
  };
\end{tikzpicture}}


\title{特殊函数简介}
\subtitle{第二讲：普西函数与泽塔函数}

\begin{document}

\begin{frame}
  \titlepage
\end{frame}

\begin{frame}{本讲提纲}
  普西函数（digamma function）是伽马函数的对数导数。有理分式的无穷级数可以收成它在极点上的值。黎曼泽塔函数（Riemann zeta function）是正整数幂次的倒数和。
  \begin{enumerate}
    \item $\psi(z)=\Gamma'(z)/\Gamma(z)$ 的递推和反射；
    \item 用部分分式求 $\sum p(n)/d(n)$；
    \item $\zeta(z)=\sum k^{-z}$，以及 $\zeta(2)=\pi^2/6$。
  \end{enumerate}
\end{frame}

\section{普西函数}

\begin{frame}{定义}
  \begin{equation*}
    \boxed{\displaystyle
      \psi(z)=\frac{\mathrm{d}}{\mathrm{d}z}\ln\Gamma(z)
      =\frac{\Gamma'(z)}{\Gamma(z)}.}
  \end{equation*}
  $\Gamma$ 在 $z=0,-1,-2,\ldots$ 有一阶极点，$\ln\Gamma$ 在那里有对数奇点，导数是一阶极点，留数都是 $-1$。除此之外 $\psi(z)$ 全平面解析。
\end{frame}

\begin{frame}{递推}
  对 $\Gamma(z+1)=z\Gamma(z)$ 取对数再求导，
  \[
    \psi(z+1)=\psi(z)+\frac{1}{z}.
  \]
  连用 $n$ 次，
  \begin{equation*}
    \psi(z+n)
    =\psi(z)+\frac{1}{z}+\frac{1}{z+1}+\cdots+\frac{1}{z+n-1}.
  \end{equation*}
  调和数就是 $H_n=\psi(n+1)-\psi(1)=\psi(n+1)+\gamma$。
\end{frame}

\begin{frame}{反射与倍乘}
  对 $\Gamma(z)\Gamma(1-z)=\pi/\sin(\pi z)$ 取对数导数，
  \begin{equation*}
    \psi(1-z)=\psi(z)+\pi\cot\pi z.
  \end{equation*}
  由此还有
  \[
    \psi(z)-\psi(-z)=-\frac{1}{z}-\pi\cot\pi z.
  \]
  倍乘公式的对数导数是
  \[
    \psi(2z)=\frac12\psi(z)+\frac12\psi\Bigl(z+\frac12\Bigr)+\ln 2.
  \]
\end{frame}

\begin{frame}{几个值}
  欧拉常数就是 $\gamma=-\psi(1)$。倍乘公式取 $z=1/2$，
  \[
    \psi(1)=\frac12\psi\Bigl(\frac12\Bigr)+\frac12\psi(1)+\ln 2,
  \]
  解出
  \[
    \psi(1)=-\gamma,
    \qquad
    \psi\Bigl(\frac12\Bigr)=-\gamma-2\ln 2.
  \]
  后面的级数要用到
  \begin{align*}
    \psi\Bigl(\frac13\Bigr)
    &=-\gamma-\frac{\pi}{2\sqrt{3}}-\frac32\ln 3,\\
    \psi\Bigl(\frac23\Bigr)
    &=-\gamma+\frac{\pi}{2\sqrt{3}}-\frac32\ln 3.
  \end{align*}
\end{frame}

\begin{frame}{大 $n$ 时像对数}
  \[
    \lim_{n\to\infty}\bigl[\psi(z+n)-\ln n\bigr]=0.
  \]
  部分分式求和时，每一项 $\psi(\alpha_k+N)$ 都换成 $\ln N$ 再相加。系数之和为零，对数就消掉，无穷级数只留下 $-\sum a_k\psi(\alpha_k)$。
\end{frame}

\section{有理级数}

\begin{frame}{部分分式}
  设 $u_n=p(n)/d(n)$，$p$ 的次数至少比 $d$ 低 $2$，级数才收敛。$d$ 的零点都是一阶时，
  \[
    u_n=\sum_{k=1}^{m}\frac{a_k}{n+\alpha_k},
    \qquad
    \sum_{k=1}^{m}a_k=0.
  \]
  递推把有限和写成普西函数，
  \[
    \sum_{n=0}^{N}u_n
    =\sum_{k=1}^{m}a_k
      \bigl[\psi(\alpha_k+N)-\psi(\alpha_k)\bigr].
  \]
  $N\to\infty$ 时，
  \begin{equation*}
    \boxed{\displaystyle
      \sum_{n=0}^{\infty}u_n
      =-\sum_{k=1}^{m}a_k\,\psi(\alpha_k).}
  \end{equation*}
\end{frame}

\begin{frame}{例：三个线性因子}
  \[
    \frac{1}{(3n+1)(3n+2)(3n+3)}
    =\frac16\frac{1}{n+1/3}
     -\frac13\frac{1}{n+2/3}
     +\frac16\frac{1}{n+1}.
  \]
  三个系数 $1/6$、$-1/3$、$1/6$ 之和是零。求和公式给出
  \[
    \sum_{n=0}^{\infty}
    =-\frac16\Biggl[
        \psi\Bigl(\frac13\Bigr)
        -2\psi\Bigl(\frac23\Bigr)
        +\psi(1)
      \Biggr].
  \]
\end{frame}

\begin{frame}{代进特殊值}
  \[
    \psi\Bigl(\frac13\Bigr)-2\psi\Bigl(\frac23\Bigr)+\psi(1)
    =-\frac{3\pi}{2\sqrt{3}}+\frac32\ln 3.
  \]
  $\gamma$ 在三项里相消。于是
  \[
    \sum_{n=0}^{\infty}
      \frac{1}{(3n+1)(3n+2)(3n+3)}
    =\frac14\Biggl(\frac{\pi}{\sqrt{3}}-\ln 3\Biggr).
  \]
\end{frame}

\begin{frame}{例：$\sum 1/(n^2+a^2)$}
  $a>0$ 时，
  \[
    \frac{1}{n^2+a^2}
    =\frac{\mathrm{i}}{2a}
      \Biggl(\frac{1}{n+\mathrm{i}a}-\frac{1}{n-\mathrm{i}a}\Biggr).
  \]
  \[
    \sum_{n=0}^{\infty}\frac{1}{n^2+a^2}
    =-\frac{\mathrm{i}}{2a}
      \bigl[\psi(\mathrm{i}a)-\psi(-\mathrm{i}a)\bigr].
  \]
  反射差给出 $\psi(\mathrm{i}a)-\psi(-\mathrm{i}a)=\mathrm{i}\bigl[1/a+\pi\coth(\pi a)\bigr]$，所以
  \begin{equation*}
    \sum_{n=0}^{\infty}\frac{1}{n^2+a^2}
    =\frac{1}{2a^2}\bigl[1+\pi a\coth(\pi a)\bigr].
  \end{equation*}
\end{frame}

\section{泽塔函数}

\begin{frame}{级数定义}
  $\operatorname{Re} z>1$ 时，
  \begin{equation*}
    \boxed{\displaystyle
      \zeta(z)=\sum_{k=1}^{\infty} k^{-z}.}
  \end{equation*}
  调和级数是 $z=1$ 的边界，那里发散。偶数值是 $\pi$ 的幂，
  \[
    \zeta(2)=\sum_{n=1}^{\infty}\frac{1}{n^2}=\frac{\pi^2}{6}.
  \]
  $\psi'(1)=\zeta(2)=\pi^2/6$。
\end{frame}

\begin{frame}{积分表示}
  \[
    \zeta(z)
    =\frac{1}{\Gamma(z)}\int_0^{\infty}\frac{t^{z-1}}{e^{t}-1}\,\mathrm{d}t,
    \qquad\operatorname{Re} z>1.
  \]
  把 $1/(e^t-1)$ 展成 $e^{-t}+e^{-2t}+\cdots$，逐项积分就是 $\Gamma(z)$ 乘 $\sum k^{-z}$。
  \vspace{0.35em}
  这个积分把 $\zeta(z)$ 从 $\operatorname{Re} z>1$ 继续延拓。$z=1$ 仍是一阶极点。
\end{frame}

\begin{frame}{这一讲留下的公式}
  \[
    \psi(z+1)=\psi(z)+\frac{1}{z},
    \qquad
    \psi(1)=-\gamma,
    \qquad
    \psi(1-z)=\psi(z)+\pi\cot\pi z.
  \]
  一阶极点的有理级数等于 $-\sum a_k\psi(\alpha_k)$。
  \[
    \zeta(2)=\frac{\pi^2}{6},
    \qquad
    \sum_{n=0}^{\infty}\frac{1}{n^2+a^2}
    =\frac{1+\pi a\coth(\pi a)}{2a^2}.
  \]
\end{frame}

\end{document}
